\section{Native determinant traces generate the actual invariant ring}
\label{sec:cubic-native-invariant-surjectivity}

The map of AC-G97 has the correct direction for native support, but
its surjectivity was not included there.  The normalization at $g_0$
now supplies integral coordinate projectors.  They recover every
weight-zero generator by a trace identity.  This also proves the
actual finite weight-fiber condition needed in the changed-tuple
dimension arguments.

Retain the actual rings $B\subset A\subset R$, representation $r$,
local characters $\zeta_{v,i}$, and native map
\[
             \Phi:\mathscr B\longrightarrow B
\tag{G102.1}
\]
of AC-G94--AC-G97.  Here $\mathscr B$ is the native complete
Noetherian local ordinary determinant ring over $\mathcal O$;
its residue field is the same finite field $k$ as that of $B$.
The characteristic of $k$ is $269$.  Write $J$ for the
antidiagonal three-by-three matrix, and use indices $0,1,2$.
The star operation on the tuple group $G$ is the one in AC-G94,
so $r(g^\star)=Jr(g)^tJ$.
The normalization element is the actual one used in AC-G98:
\[
 a_0=\theta_0(g_0)-\theta_2(g_0)\ne0.
\tag{G102.2}
\]
In particular its residual endpoints are inverse and neither
has square one.  This is a retained separation
condition, not a new assertion about every possible $g_0$.

\begin{theorem}[AC-G102-surj: integral trace generation for the cubic tuple]
\label{thm:ACG102surj}
The map (G102.1) is surjective.  More precisely, its image
contains the local characters and every zero-weight monomial
in the entries of $r$.  The proof uses only the inverses of
$2$ and one element whose residue is $a_0$; it does not invert
the coefficient prime or assume that any ring is reduced.
\end{theorem}

\begin{proof}
\emph{1. The image is a closed complete local ring.}
Put $C=\Phi(\mathscr B)\subset B$.  A complete Noetherian
local ring with finite residue field is compact: it is the
inverse limit of its finite quotients by powers of its maximal
ideal.  The map $\Phi$ is continuous by AC-G97, and $B$ is
Hausdorff.  Thus $C$ is compact and closed in $B$.
Its kernel in $\mathscr B$ is an ideal in a complete Noetherian
local ring and is closed.  Consequently
$C\simeq\mathscr B/\ker\Phi$ is complete Noetherian local,
with its quotient topology.  The continuous bijection from
this compact quotient to its image in $B$ is a homeomorphism.
Its residue field is $k$: the map is local, and its image
contains $\mathcal O$, which maps onto $k$.

For clarity the maximal ideal of $C$ is $C\cap\mathfrak m_B$.
If $c=\Phi(b)$ has nonzero residue, then $b$ is a unit in
$\mathscr B$, so its inverse maps to an inverse in $C$.
The other elements lie in the image of the maximal ideal.
This proves the assertion without identifying the ambient
and quotient adic topologies by fiat.

Every trace $\operatorname{tr}r(g)$ belongs to $C$.
Choose a lift $\widetilde g$ of $g$ under the surjection of
AC-G97-link.  Its native representation is
$\widetilde r(\widetilde g)=\kappa(\widetilde g)r(g)$,
where $\kappa$ is the fixed scalar character with values in
$\mathcal O^\times$.  The negative coefficient of degree two of its
degree-three characteristic polynomial gives its trace,
which is the image of the native determinant-law trace.
Multiplication by $\kappa(\widetilde g)^{-1}$ gives
$\operatorname{tr}r(g)\in C$.
The same scalar untwist of the native local characters gives
\[
                   \zeta_{v,i}(h)\in C
                   \quad\text{for every }v,i,h.
\tag{G102.3}
\]
Only the already proved Pair map is used; no surjectivity
has been presumed in this step.

\emph{2. Recover the normalization scalar from traces.}
Let
\[
                         Q=r(g_0)-r(g_0^\star).
\tag{G102.4}
\]
The normalization makes $Q$ diagonal.  Polarity gives
$Q=-JQ^tJ$, so, since $2$ is a unit,
\[
                 Q=\operatorname{diag}(a,0,-a),
                 \qquad \overline a=a_0.
\tag{G102.5}
\]
In particular $a$ is a unit in $R$ and belongs to $B$ as a
diagonal-entry difference.  Expanding $Q^2$ and using
$r(g)r(h)=r(gh)$ expresses its trace as an integral linear
combination of traces of group elements.  Hence
\[
                         b=\tfrac12\operatorname{tr}(Q^2)
                           =a^2\in C.
\tag{G102.6}
\]
The polynomial $X^2-b\in C[X]$ has the simple residual root
$a_0$, with derivative $2a_0\ne0$.  Complete-local Hensel
lifting therefore supplies $a_C\in C$ with $a_C^2=b$ and
residue $a_0$.  The precise standard input is Stacks,
Lemma~10.153.9 in tag~04GE, whose whole proof was read.
In $R$ the equality
$(a-a_C)(a+a_C)=0$ and the unit residue $2a_0$ of $a+a_C$
give $a=a_C$.  Thus $a,a^{-1}\in C$.
This also proves uniqueness in the potentially nonreduced
ambient ring, without using a fraction field.

The three coordinate projectors are consequently polynomials
in $Q$ with coefficients in $C$:
\[
 E_0=\frac{Q^2+aQ}{2a^2},\qquad
 E_1=1-\frac{Q^2}{a^2},\qquad
 E_2=\frac{Q^2-aQ}{2a^2}.
\tag{G102.7}
\]
Here $E_i$ is the diagonal matrix with a single one in
position $i$.  Formula (G102.5) verifies (G102.7) directly.

\emph{3. Every directed cycle of entries is a native trace.}
Let $i_0,i_1,\ldots,i_s=i_0$ be indices and let
$g_1,\ldots,g_s\in G$.  Ordinary matrix multiplication gives
\[
 \prod_{j=1}^{s}r_{i_{j-1},i_j}(g_j)
 =\operatorname{tr}
   \bigl(E_{i_0}r(g_1)E_{i_1}r(g_2)\cdots
                         E_{i_{s-1}}r(g_s)E_{i_0}\bigr).
\tag{G102.8}
\]
Substitute (G102.7), expand (G102.4), and multiply out.
The right side becomes a finite $C$-linear combination of
traces of products of group matrices.  Each such product
is a group matrix, so Step~1 places its trace in $C$.
Thus every cycle (G102.8) lies in $C$.
For $s=1$ this includes every diagonal entry.

It remains to account for the fact that our grading uses
the one-dimensional polarized torus, rather than the full
diagonal torus.  The exact entry identities from polarity are
\[
 r_{01}(g^\star)=r_{12}(g),\qquad
 r_{10}(g^\star)=r_{21}(g),\qquad
 r_{02}(g^\star)=r_{02}(g),\qquad
 r_{20}(g^\star)=r_{20}(g).
\tag{G102.9}
\]
The weights of the first two adjacent positions are $1$,
the opposite adjacent positions have weight $-1$, and
positions $02,20$ have weights $2,-2$.
By (G102.9), either adjacent entry can be placed at either
position of the same weight, after replacing its group
argument by its star when necessary.

The complete list of nontrivial zero-weight atoms in
AC-G95 can therefore be realized as the following cycles:
\[
\begin{array}{c|c}
 \text{weight pattern}&\text{entry product after (G102.9)}\\ \hline
 (1,-1)&r_{01}(g)r_{10}(h)\\
 (2,-2)&r_{02}(g)r_{20}(h)\\
 (1,1,-2)&r_{01}(g)r_{12}(h)r_{20}(u)\\
 (2,-1,-1)&r_{02}(g)r_{21}(h)r_{10}(u).
\end{array}
\tag{G102.10}
\]
Each belongs to $C$ by (G102.8).  Factors have commutative
scalar values in $R$, so they may be reordered to these
patterns.  Removing a weight-zero variable or one of the
four atoms successively decomposes every zero-weight
monomial; the elementary exhaustive argument is the one
proved in AC-G95.  Thus every such monomial belongs to $C$.

By (G102.3), $C$ also contains every local-character generator
of $B$.  It contains $\mathcal O$.  Finally $C$ is closed in
$R$: it is closed in $B$, and $B$ is the closed subring of
$R$ constructed in AC-G95.  The defining closure of these
generators is exactly $B$.  Hence $B\subset C$, and $C=B$.
This proves surjectivity.
\end{proof}

\paragraph{The actual weight-fiber ideal.}
For a coefficient-prime local inertia group $I_v$, put
\[
 K_v=I_v\cap\bigcap_{i=0}^{2}
                  \ker(\theta_i\circ\operatorname{loc}_v),
 \qquad
 \lambda_\theta=(\varpi)B+
  (\,\zeta_{v,i}(g)-1:g\in K_v,\ v,i\,)B.
\tag{G102.11}
\]
Here $\varpi$ is a uniformizer of $\mathcal O$.
The quotient $I_v/K_v$ is finite of order prime to $p$:
it embeds into the finite group $(k^\times)^3$.
Every generator in (G102.11) has zero residue, so
$\lambda_\theta\subset\mathfrak m_B$.
This is the changed-tuple form of the source's
\path{Space.lamId}.  Its kernel must be the joint residual
kernel.  Using only the kernel of the rational character
could include elements on which the cubic character is
nontrivial and would give a unit among the proposed generators.

\begin{theorem}[AC-G102-finlam: actual finite weight fiber from native support]
\label{thm:ACG102finlam}
Let $J_{\rm aut}\subset\mathscr B$ be the proper native ideal
of AC-G89 on the current host, with the retained source
\path{ModU}, \path{LamEx} and support hypotheses used there
and in AC-G93-fill.  Let $\Lambda$ be its actual native
\path{LamObj}.  Put
\[
                        I=\Phi(J_{\rm aut})B.
\tag{G102.12}
\]
Then $I$ is proper, $B/I$ is finite over $\Lambda$, and
\[
 \dim_k B/(I+\lambda_\theta)<\infty,
 \qquad \sqrt{I+\lambda_\theta}=\mathfrak m_B.
\tag{G102.13}
\]
In particular the exact tuple version of the source
\path{FinLam} condition holds.
No torsionfreeness of $B/I$ over $\Lambda$, persistence of
all native components, or support of a prescribed
varying-character curve is asserted.
\end{theorem}

\begin{proof}
Surjectivity gives a surjection
\[
             \mathscr B/J_{\rm aut}\longrightarrow B/I.
\tag{G102.14}
\]
The native source \path{RingProps.fin}, retained exactly
in AC-G89, says that the left side is finite over $\Lambda$.
Its quotient is therefore finite as well.
The local map $\Phi$ takes $J_{\rm aut}$ into
$\mathfrak m_B$, so $I$ is proper.

Set $D=B/(I+\lambda_\theta)$.  This is a nonzero complete
Noetherian local $k$-algebra with residue field $k$.
It is an allowed native test object over $\mathcal O$;
we have killed $\varpi$, not merely a power of $p$.
Write the image of a local character in $D^\times$ as
$z_{v,i}$.  Formula (G102.11) makes it trivial on $K_v$.
It factors through $I_v/K_v$ on inertia.
For $g\in I_v$, let $m$ be the order of its image in that
finite prime-to-$p$ quotient.  Both $z_{v,i}(g)$ and its
residual value $\theta_i(\operatorname{loc}_v g)\in k^\times$
have $m$th power one.  Their quotient $u$ has residue one and
$u^m=1$.  But
\[
 (u-1)(1+u+\cdots+u^{m-1})=0,
\tag{G102.15}
\]
and the second factor has nonzero residue $m$.
It is a unit in the local ring $D$, so $u=1$.
Consequently every $z_{v,i}|_{I_v}$ is exactly its fixed
residual character.  This argument permits nilpotents.

The native local character is the product of this tuple
character with the fixed scalar $\kappa$ of AC-G97.
The latter has values in $\mathcal O^\times$ and, in $D$,
is exactly its residue.  Thus every native inertial character
over $D$ is its fixed residual character, including its
finite scalar factor.  No assumption that the full
cyclotomic character is trivial on $K_v$ is needed.

The source universal property \path{LamObj.univ} now
identifies the map $\Lambda\to D$ with
\[
                           \Lambda\longrightarrow k
                                      \longrightarrow D.
\tag{G102.16}
\]
Indeed both maps induce precisely those same ordered
inertial characters.  For the second map use the residue
of the universal native inertia datum, followed by the
coefficient embedding $k\hookrightarrow D$; its extension
to each full local group is the residue of the corresponding
native local character.  It is therefore an admissible
inertial datum.  The uniqueness clause of the actual
\path{LamObj} applies, and AC-G97.9 identifies the first map
with the native weight map through $B$.

As a quotient of the finite $\Lambda$-module $B/I$, the
ring $D$ has finitely many module generators.  Since its
$\Lambda$-action factors through $k$ by (G102.16), these
same generators span $D$ over $k$.
This proves the first assertion of (G102.13).
A finite-dimensional local $k$-algebra has nilpotent
maximal ideal: the descending powers stabilize, and
Nakayama applied to the stabilized finite module makes
that power zero.  Hence some power of $\mathfrak m_B$
is contained in $I+\lambda_\theta$.
The reverse radical inclusion follows from
$I+\lambda_\theta\subset\mathfrak m_B$.
This proves (G102.13) and the stated \path{FinLam} condition.
\end{proof}

\paragraph{Exact sources and geometric interpretation.}
The complete native \path{RepAlg.lean} and
\path{UniversalOrdinaryDeterminantRing.lean} were freshly
read, including the polynomial-law trace, Pair maps,
\path{LamObj.univ} and \path{LamObj.toPair_spec}.
The exact \path{RingProps.fin}, \path{ModU}, \path{LamEx}
and \path{ProModularRingExists} declarations are in
\path{WeilGroupCharacters/ArithmeticPointClassicality.lean}.
The source \path{Space.lamId} and \path{FinLam} are in
\path{FrobeniusSmoothVectors/LeAdicClosure.lean};
\path{Setup.PhiEx} and \path{Setup.FinLamOf} are in
\path{CurveDVRSpecialization/SecondThirdPrimeConditionTransferAlt.lean}.
These files are at the previously pinned commit
\texttt{7cbf7a0c32fca1badf05931879a37c3f8a2b03b4}.
The new argument proves the actual tuple transfer and does
not instantiate the old one-character \path{TwLink}.
The complete applied Hensel proof is Stacks,
Lemma~10.153.9, with uniqueness Lemma~10.153.2,
\url{https://stacks.math.columbia.edu/tag/04GE}.
All upstream automorphic support propositions retain their
earlier explicit status.  No Lean execution is claimed.

The surjection (G102.1) realizes the formal invariant tuple
space as a closed formal subspace of the native determinant
space.  Equation (G102.14) pulls native support to that
subspace, and (G102.13) proves that its fixed residual
weight fiber is zero dimensional.  These are concrete
complete-local algebra and determinant/local-system
statements in the existing SGA~1 and Stacks dictionary.
No new Grothendieck archive scan is attributed.
The actual adjacent finite upper certificate, remaining
local presentation and propagation arguments, original
counting, and prescribed varying-character curve remain
separate from this finite weight-fiber theorem.
